% Problem dossier: VI.03.P2
\subsection*{Problem VI.03.P2 --- Polynomial interpolation on a finite affine set}
\label{prob:vi03-p2}
\paragraph{Problem.}
Let \(X=\{a_1,\dots,a_r\}\subseteq\mathbb A_k^n\) be distinct rational points.  Prove that
the evaluation map
\[
k[x_1,\dots,x_n]\to k^r,
\qquad f\mapsto(f(a_1),\dots,f(a_r))
\]
is surjective and deduce \(k[X]\cong k^r\).
\ifdefined\FullProblemDossiers
\paragraph{Why this problem matters.}
Finite algebraic sets behave like completely independent collections of points: polynomial
functions can prescribe arbitrary values at each point.
\paragraph{Useful ideas before starting.}
For each pair \(i\neq j\), choose a coordinate on which \(a_i\) and \(a_j\) differ.  Build a
linear factor that is \(1\) at \(a_i\) and \(0\) at \(a_j\), then multiply these factors.
\paragraph{Guided hints.}
\textbf{Hint 1.} Construct \(e_i\) with \(e_i(a_i)=1\) and \(e_i(a_j)=0\) for \(j\neq i\).
\textbf{Hint 2.} Then \(\sum c_i e_i\) realizes any prescribed values \(c_i\).
\paragraph{Solution.}
For fixed \(i\) and each \(j\neq i\), choose a coordinate \(\ell(i,j)\) with
\((a_i)_{\ell(i,j)}\neq(a_j)_{\ell(i,j)}\), and set
\[
L_{ij}(x)=\frac{x_{\ell(i,j)}-(a_j)_{\ell(i,j)}}{(a_i)_{\ell(i,j)}-(a_j)_{\ell(i,j)}}.
\]
Then \(L_{ij}(a_i)=1\) and \(L_{ij}(a_j)=0\).  Put
\[
e_i=\prod_{j\neq i}L_{ij}.
\]
Thus \(e_i(a_i)=1\) and \(e_i(a_j)=0\) for \(j\neq i\).  Given
\((c_1,\dots,c_r)\in k^r\), the polynomial \(\sum_i c_i e_i\) has exactly these values, so
evaluation is surjective.  Its kernel is \(I(X)\), hence the first isomorphism theorem gives
\(k[X]\cong k^r\).
\paragraph{What happened mathematically.}
The functions \(e_i\) are orthogonal idempotent coordinate selectors after passing to
\(k[X]\).  They separate the finite set into its individual points.
\paragraph{Extensions.}
Show \(e_i^2=e_i\), \(e_ie_j=0\) for \(i\neq j\), and \(\sum_i e_i=1\) in \(k[X]\).
\paragraph{Provenance.} New synthesis problem; it consolidates the finite-point material that was previously scattered across routine exercises.
\fi
